\chapter{Conclusiones y Trabajo Futuro}
\label{cap:conclusiones}

En este capítulo final se sintetizan las principales aportaciones científicas y técnicas desarrolladas a lo largo de este Trabajo de Fin de Máster, se discuten las limitaciones asumidas durante el modelado y se plantean las líneas prioritarias de investigación futura.

\section{Conclusiones Principales}

El presente trabajo ha desarrollado e integrado con éxito un framework completo y unificado para la optimización de operaciones de búsqueda y rescate (SAR) mediante vehículos aéreos no tripulados (UAVs) sobre escenarios topográficos reales en entornos de incertidumbre:

\begin{enumerate}
    \item \textbf{Integración Topológica y Probabilística Realista:} Se ha superado la dependencia tradicional de mapas sintéticos simplificados mediante la conexión directa con OpenStreetMap (OSM) y la parametrización de perfiles canónicos del manual de comportamiento de personas desaparecidas (\textit{Lost Person Behavior} - LPB de Robert Koester).
    \item \textbf{Filtrado de Restricciones Físicas de Transitabilidad:} La incorporación de máscaras espaciales jerárquicas con árboles R-Tree ha permitido anular estrictamente la probabilidad ($P=0.0$) en zonas no transitables (lagos y edificaciones cerradas), eliminando falsos positivos en la asignación de recursos aéreos.
    \item \textbf{Arquitectura Modular y Portabilidad (`sarenv-mts`):} Se ha establecido una pasarela middleware que transforma las coordenadas globales UTM ($\text{EPSG:32630}$) en matrices discretas operables por el motor de búsqueda MTS, junto con un evaluador estandarizado de métricas SAR (\texttt{PathEvaluatorTFM}).
    \item \textbf{Calibración Bayesiana y Eficacia de Algoritmos Bioinspirados:} A través de un exhaustivo benchmark de 900 simulaciones de Montecarlo calibradas con Optuna, se ha demostrado cuantitativamente que los planificadores metaheurísticos bioinspirados (en particular la Colonia de Abejas Artificiales - ABC y el Algoritmo de Agujeros Negros - BHA) superan sistemáticamente a los patrones de barrido geométrico ciego (Lawnmower) y a la búsqueda voraz miope, logrando una reducción significativa en los tiempos de respuesta y una mayor concentración en áreas de alta verosimilitud.
\end{enumerate}

\section{Limitaciones del Trabajo}

A pesar de la consistencia metodológica del framework desarrollado, es necesario explicitar las siguientes limitaciones asumidas en la investigación:

\begin{enumerate}
    \item \textbf{Ausencia de Penalización Física en el Vector de Control Cinemático:} 
    Las restricciones físicas de no transitabilidad (masas de agua y edificaciones) se han aplicado como una \textit{restricción dura de probabilidad inicial} ($b(v^0) = 0.0$), garantizando que los planificadores no concentren esfuerzos de búsqueda en dichas zonas. Sin embargo, no se ha incorporado un término de repulsión cinemática en las acciones de vuelo del UAV, permitiendo que el dron pueda sobrevolar en línea recta masas de agua durante su desplazamiento hacia otras regiones de alta probabilidad.

    \item \textbf{Asunción de Víctimas Estáticas en la Batería Experimental:} 
    Las 900 simulaciones del benchmark consideran que la persona desaparecida permanece en una ubicación fija e inalterable desde el inicio de la misión hasta su hallazgo o el agotamiento de la autonomía de vuelo.
\end{enumerate}

\section{Líneas de Trabajo Futuro}

Como extensión directa de los resultados y la arquitectura desarrollada, se plantean las siguientes líneas de trabajo futuro:

\begin{itemize}
    \item \textbf{Integración Operativa Aire-Tierra y Validación con Simulacros Reales de Bomberos:} 
    El escenario seleccionado en este trabajo (Casa de Campo de Madrid) responde a un ejercicio operativo real de búsqueda y salvamento coordinado por cuerpos de bomberos y protección civil, del cual se dispone de la sectorización territorial, la ubicación del Puesto de Mando Avanzado (PMA) y los registros de posicionamiento GPS (\textit{tracks} GPX) de 12 equipos de batida terrestre. Como continuación directa, se plantea articular una arquitectura colaborativa multinivel que combine dos escalas operativas:
    \begin{itemize}
        \item \textit{Escala Macro (Exploración Aérea Temprana):} El UAV opera a $10\text{ m/s}$ ($36\text{ km/h}$) a $50\text{ m}$ de altitud como vector táctico de detección rápida sobre el macro-escenario completo. Su cometido prioritario es absorber la mayor masa de creencia bayesiana en los primeros 15--20 minutos, transmitiendo las coordenadas del contacto o las zonas de máxima certidumbre al PMA para focalizar el despliegue de los recursos a pie.
        \item \textit{Escala Micro (Validación y Batida Sectorizada):} Restringir los algoritmos bioinspirados de \texttt{sarenv-mts} al polígono específico de uno de los subsectores asignados a los bomberos para contrastar cuantitativamente la trayectoria y el tiempo de cobertura del dron frente a la traza GPX registrada a pie por el equipo terrestre ($3\text{--}4\text{ km/h}$). Asimismo, cruzar el mapa de creencia residual $b(v^k)$ con las áreas batidas por los rescatadores a pie permitirá optimizar la reasignación dinámica de cuadrantes y eliminar búsquedas redundantes.
    \end{itemize}
    
    \item \textbf{Modelado de Víctimas Dinámicas mediante Cadenas de Markov y Random Walk:} 
    Incorporar matrices de transición dinámicas $P(x, y, t)$ orientadas por la pendiente del terreno, senderos y perfiles psicológicos de dispersión extraídos del Manual de Búsqueda y Salvamento Terrestre \cite{romero2018manual}.
    
    \item \textbf{Control de Vuelo Guiado por Campos Potenciales de Repulsión:} 
    Implementar restricciones cinemáticas de no sobrevuelo en la capa de bajo nivel del UAV para impedir físicamente que el dron cruce masas de agua o zonas de exclusión aérea.
\end{itemize}
